\subsection{三高・四高に共通する正値性の短い判定}
\label{nf24:positive-graph}

四rayの全 pair products が狭義正のとき，正の weights が存在することは，
正符号グラフが孤立頂点を持たないことと同値である．
同値な符号候補は $2K_2$ と $P_4$ の二つである．
三高では既存の整数障害が $P_4$ も除くが，四高にその排除を流用しない．

再利用する恒等式は次である．孤立頂点を $h$，$c_i=-p_{hi}>0$ と置くと
\[
 g_{ij}=q_i c_j-c_i p_{ij},\qquad
 \sum_{j\ne i,h}g_{ij}=1.
\]
二つの整数の和が一なので，それぞれの頂点から $g_{ij}\le0$ の辺を選べる．
選んだ正の閉路の gain は一以上で，正 weights に反する．
この議論は整数性を必要とする．零 pair のある二高以下には適用しない．

四高の完全マッチング例は既存例を用いればよい．小さい分母の局所列挙で
$P_4$ が見つからなかったことは，全尺度の排除証明ではない．
今後の四高研究では $P_4$ の独立な排除または構成と，完全マッチング型の
大域 clutching を分ける．どちらの候補でも正値性を改めて実数最適化する必要はなく，
木の各辺で開区間内の有理 weight 比を選べる．
